% AA-57
\chapter{Algebry Liego. Twierdzenia Engela i Liego}
\label{aa:ch:lie-algebras}

W rozdziale~\ref{aa:ch:root-systems} rozważaliśmy układy
korzeni i grupy Weyla. Do algebraicznej interpretacji tych
obiektów potrzebujemy teorii algebr Liego. Zaczniemy od
rachunku komutatorów, a następnie wykażemy dwa podstawowe
twierdzenia o reprezentacjach algebr nilpotentnych
i rozwiązalnych.

\section{Nawias Liego i przykłady macierzowe}

\begin{definicja}[Algebra Liego]
\label{aa:def:lie-algebra}
Niech $K$ będzie ciałem. Algebrą Liego nad $K$ nazywamy
przestrzeń wektorową $\mathfrak g$ nad $K$ z dwuliniowym
odwzorowaniem $[\cdot,\cdot]:\mathfrak g\times\mathfrak g
\to\mathfrak g$, zwanym nawiasem Liego, spełniającym
dla wszystkich $x,y,z\in\mathfrak g$:
\[
 [x,x]=0,\qquad
 [x,[y,z]]+[y,[z,x]]+[z,[x,y]]=0.
\]
Drugą równość nazywamy tożsamością Jacobiego.
Podalgebrą Liego jest podprzestrzeń zamknięta
na nawias. Dla podprzestrzeni $U,W\subset\mathfrak g$
przez $[U,W]$ rozumiemy rozpiętość liniową wszystkich
$[u,w]$, gdzie $u\in U,w\in W$.
\end{definicja}

Z rozwinięcia $[x+y,x+y]=0$ wynika
$[x,y]=-[y,x]$, także w charakterystyce dwa.
Tożsamość Jacobiego można zapisać w postaci
\[
 [x,[y,z]]=[[x,y],z]+[y,[x,z]].
\]
Istotnie przenosimy dwa pozostałe składniki
na prawą stronę i używamy antysymetrii.

\begin{przyklad}[Algebra łączna z komutatorem]
\label{aa:prz:associative-lie}
Niech $A$ będzie algebrą łączną nad $K$.
Na jej przestrzeni wektorowej połóżmy $[x,y]=xy-yx$.
Nawias jest dwuliniowy i $[x,x]=0$.
Ponadto
\[
\begin{aligned}
 [x,[y,z]]&=xyz-xzy-yzx+zyx,\\
 [y,[z,x]]&=yzx-yxz-zxy+xzy,\\
 [z,[x,y]]&=zxy-zyx-xyz+yxz.
\end{aligned}
\]
Suma trzech prawych stron znika wyraz po wyrazie.
Otrzymujemy algebrę Liego. Dla $A=\operatorname{End}_K(V)$
oznaczamy ją przez $\mathfrak{gl}(V)$, a dla
$V=K^n$ przez $\mathfrak{gl}_n(K)$.
Dla jednostek macierzowych $E_{ij}$ zachodzi
$E_{ij}E_{kl}=\delta_{jk}E_{il}$, ponieważ obie
strony mają to samo działanie na wektorach standardowych.
Zatem
\[
 [E_{ij},E_{kl}]=\delta_{jk}E_{il}-\delta_{li}E_{kj}.
\]
\end{przyklad}

\begin{przyklad}[Algebry klasyczne]
\label{aa:prz:classical-lie}
Niech $K$ będzie ciałem charakterystyki zero
i $n\ge1$. Określamy
\[
\begin{aligned}
 \mathfrak{sl}_n(K)&=\{X\in M_n(K):\operatorname{tr}X=0\},\\
 \mathfrak{so}_n(K)&=\{X\in M_n(K):X^{\mathsf T}+X=0\},\\
 \mathfrak{sp}_{2n}(K)&=\{X\in M_{2n}(K):X^{\mathsf T}J+JX=0\},
 \qquad J=\begin{pmatrix}0&I_n\\-I_n&0\end{pmatrix}.
\end{aligned}
\]
Wszystkie są podprzestrzeniami. Ponieważ
$\operatorname{tr}(XY)=\sum_{i,j}x_{ij}y_{ji}
=\operatorname{tr}(YX)$, każdy komutator ma ślad zero,
więc $\mathfrak{sl}_n$ jest podalgebrą, a nawet
ideałem w $\mathfrak{gl}_n$ w sensie poniższej definicji.

Jeżeli dla ustalonej macierzy $B$ mamy
$X^{\mathsf T}B=-BX$, $Y^{\mathsf T}B=-BY$, to
\[
 [X,Y]^{\mathsf T}B
 =Y^{\mathsf T}X^{\mathsf T}B-X^{\mathsf T}Y^{\mathsf T}B
 =BYX-BXY=-B[X,Y].
\]
Podstawienie $B=I_n$ lub $B=J$ dowodzi zamkniętości
pozostałych dwóch przestrzeni na komutator.

Wymiary wynoszą odpowiednio $n^2-1$, $n(n-1)/2$
oraz $n(2n+1)$. Pierwsza przestrzeń jest jądrem
niezerowego funkcjonału śladu (wartość w $E_{11}$
wynosi $1$). W drugiej wpisy nad przekątną są dowolne,
pod przekątną są ich przeciwieństwami, a przekątna
jest zerowa, ponieważ $2\ne0$ w $K$.
W trzeciej piszemy $X=\begin{pmatrix}P&Q\\R&S\end{pmatrix}$.
Mnożenie bloków daje
\[
 X^{\mathsf T}J+JX
 =\begin{pmatrix}R-R^{\mathsf T}&P^{\mathsf T}+S\\
 -S^{\mathsf T}-P&Q^{\mathsf T}-Q\end{pmatrix}.
\]
Zatem $S=-P^{\mathsf T}$, a $Q,R$ są symetryczne.
Macierz $P$ ma $n^2$ dowolnych współczynników,
$Q,R$ po $n(n+1)/2$, skąd wskazany wymiar.
Przy $n=1$ mamy $\mathfrak{sl}_1=\mathfrak{so}_1=0$
i $\mathfrak{sp}_2=\mathfrak{sl}_2$.
\end{przyklad}

\section{Ideały, ilorazy i reprezentacje}

\begin{definicja}[Ideał i homomorfizm]
\label{aa:def:lie-ideal}
Ideałem algebry Liego $\mathfrak g$ nazywamy
podprzestrzeń $\mathfrak a\subset\mathfrak g$ taką,
że $[\mathfrak g,\mathfrak a]\subset\mathfrak a$.
Homomorfizm algebr Liego nad $K$ jest liniowym
odwzorowaniem $f:\mathfrak g\to\mathfrak h$ spełniającym
$f([x,y])=[f(x),f(y)]$ dla wszystkich $x,y$.
\end{definicja}

\begin{twierdzenie}[Algebra ilorazowa i twierdzenie o izomorfizmie]
\label{aa:tw:lie-quotient}
Dla ideału $\mathfrak a$ algebry Liego $\mathfrak g$
nawias $[x+\mathfrak a,y+\mathfrak a]=[x,y]+\mathfrak a$
określa algebrę Liego na przestrzeni ilorazowej.
Jeśli $f:\mathfrak g\to\mathfrak h$ jest homomorfizmem,
to $\ker f$ jest ideałem, $\operatorname{im}f$
podalgebrą, a odwzorowanie indukowane
$\mathfrak g/\ker f\to\operatorname{im}f$ jest
izomorfizmem algebr Liego.
\end{twierdzenie}
\begin{proof}
Zmiana reprezentantów o $u,v\in\mathfrak a$ zmienia
nawias o $[x,v]+[u,y]+[u,v]\in\mathfrak a$.
Nawias ilorazowy jest więc dobrze określony.
Dwuliniowość, alternowanie i tożsamość Jacobiego
wynikają z odpowiednich własności w $\mathfrak g$
po zastosowaniu rzutu ilorazowego.
Dla $u\in\ker f$ mamy $f([x,u])=[f(x),0]=0$,
a nawias dwóch elementów obrazu jest obrazem nawiasu.
Odwzorowanie $x+\ker f\mapsto f(x)$ jest niezależne
od reprezentanta, surjektywne na obraz i ma jądro
zerowe. Zachowuje działania liniowe i nawias
na mocy definicji, więc jest żądanym izomorfizmem.
\end{proof}

\begin{definicja}[Reprezentacja i moduł]
\label{aa:def:lie-representation}
Reprezentacją algebry Liego $\mathfrak g$ nad $K$
w przestrzeni wektorowej $V$ nad $K$ nazywamy
homomorfizm $\rho:\mathfrak g\to\mathfrak{gl}(V)$.
Równoważnie określamy dwuliniowe działanie
$x\cdot v=\rho(x)v$, spełniające
$[x,y]\cdot v=x\cdot(y\cdot v)-y\cdot(x\cdot v)$.
Przestrzeń $V$ z takim działaniem nazywamy
$\mathfrak g$-modułem. Podmoduł to podprzestrzeń
$U\subset V$ zachowana przez wszystkie $\rho(x)$.
Reprezentacja jest wierna, gdy $\ker\rho=0$.
\end{definicja}

Dla podmodułu $U$ działanie na $V/U$ określamy przez
$x\cdot(v+U)=x\cdot v+U$. Zmiana $v$ o element $U$
zmienia wynik o element $U$, więc działanie jest dobrze
określone. Równość definiującą moduł otrzymujemy
przez przejście do klas w odpowiedniej równości w $V$.
Reprezentacja naturalna dowolnej podalgebry
$\mathfrak{gl}(V)$ jest jej włożeniem w tę algebrę,
a zatem jest wierna.

\begin{definicja}[Centrum i reprezentacja dołączona]
\label{aa:def:adjoint-lie}
Dla algebry Liego $\mathfrak g$ definiujemy
$Z(\mathfrak g)=\{x\in\mathfrak g:[x,y]=0
\text{ dla każdego }y\in\mathfrak g\}$ oraz
$\operatorname{ad}_x(y)=[x,y]$ dla $x,y\in\mathfrak g$.
\end{definicja}

\begin{stwierdzenie}
\label{aa:prop:adjoint-lie}
Dla dowolnej algebry Liego $\mathfrak g$ odwzorowanie
$\operatorname{ad}:\mathfrak g\to\mathfrak{gl}(\mathfrak g)$
jest reprezentacją o jądrze $Z(\mathfrak g)$.
\end{stwierdzenie}
\begin{proof}
Liniowość wynika z dwuliniowości nawiasu.
Dla $x,y,z\in\mathfrak g$ tożsamość Jacobiego daje
\[
 (\operatorname{ad}_x\operatorname{ad}_y-
 \operatorname{ad}_y\operatorname{ad}_x)(z)
 =[x,[y,z]]-[y,[x,z]]=[[x,y],z].
\]
Zatem $\operatorname{ad}_{[x,y]}=
[\operatorname{ad}_x,\operatorname{ad}_y]$.
Warunek $\operatorname{ad}_x=0$ jest dokładnie
warunkiem przynależności do centrum.
\end{proof}

\section{Derywacje i automorfizmy}

\begin{definicja}[Derywacja i automorfizm]
\label{aa:def:lie-derivation}
Derywacją algebry Liego $\mathfrak g$ nazywamy
liniowe odwzorowanie $D:\mathfrak g\to\mathfrak g$
spełniające $D[x,y]=[Dx,y]+[x,Dy]$ dla wszystkich
$x,y\in\mathfrak g$. Zbiór derywacji oznaczamy
$\operatorname{Der}(\mathfrak g)$.
Automorfizmem nazywamy bijektywny homomorfizm
$\mathfrak g\to\mathfrak g$.
\end{definicja}

\begin{stwierdzenie}
\label{aa:prop:lie-derivations}
Dla dowolnej algebry Liego $\mathfrak g$ zbiór
$\operatorname{Der}(\mathfrak g)$ jest podalgebrą
$\mathfrak{gl}(\mathfrak g)$, a
$\operatorname{ad}(\mathfrak g)$ jest w nim ideałem.
\end{stwierdzenie}
\begin{proof}
Warunek derywacji jest liniowy względem operatora.
Dla derywacji $D,E$ rozwijamy
\[
\begin{aligned}
 DE[x,y]&=[DEx,y]+[Ex,Dy]+[Dx,Ey]+[x,DEy],\\
 ED[x,y]&=[EDx,y]+[Dx,Ey]+[Ex,Dy]+[x,EDy].
\end{aligned}
\]
Odejmując, otrzymujemy warunek derywacji dla $[D,E]$.
Tożsamość Jacobiego mówi, że $\operatorname{ad}_x$
jest derywacją. Dla dowolnej derywacji $D$ zachodzi
\[
 [D,\operatorname{ad}_x](y)=D[x,y]-[x,Dy]=[Dx,y].
\]
Zatem $[D,\operatorname{ad}_x]=\operatorname{ad}_{Dx}$,
co dowodzi własności ideału.
\end{proof}

\begin{stwierdzenie}[Eksponenta derywacji]
\label{aa:prop:exp-derivation}
Niech $\mathfrak g$ będzie skończenie wymiarową
algebrą Liego nad $\mathbb C$ lub $\mathbb R$.
Jeśli $D$ jest jej derywacją, to $\exp(tD)$ jest
automorfizmem dla każdego $t\in\mathbb R$.
\end{stwierdzenie}
\begin{proof}
W skończonym wymiarze szereg
$\sum_{k\ge0}t^kD^k/k!$ zbiega w normie operatorowej,
bo normy wyrazów są ograniczone przez
$|t|^k\|D\|^k/k!$. Mnożenie dwóch bezwzględnie
zbieżnych szeregów dla $t,-t$ daje $I$:
w stopniu $k\ge1$ współczynnik wynosi
$\sum_{j=0}^k(-1)^j/(j!(k-j)!)=0$.
Zatem $\exp(tD)$ jest odwracalne.
Indukcja z użyciem własności derywacji i tożsamości
Pascala daje dla każdego $k\ge0$
\[
 D^k[x,y]=\sum_{j=0}^k\binom{k}{j}[D^jx,D^{k-j}y].
\]
Dla $k=0$ jest to równość identyczna. Po zastosowaniu
$D$ każdy składnik daje dwa składniki, których
współczynniki w stopniu $k+1$ sumują się według
tożsamości Pascala, także na końcach.
Nawias jako odwzorowanie dwuliniowe w skończonym
wymiarze jest ciągły i spełnia
$\|[u,v]\|\le c\|u\|\|v\|$ dla pewnego $c$.
Szereg podwójny powstały po rozwinięciu
$[\exp(tD)x,\exp(tD)y]$ jest więc bezwzględnie
zbieżny (majorantą jest $c\|x\|\|y\|e^{2|t|\|D\|}$).
Grupując wyrazy według sumy indeksów i korzystając
z wyprowadzonego wzoru, otrzymujemy
$[\exp(tD)x,\exp(tD)y]=\exp(tD)[x,y]$.
\end{proof}

\section{Nilpotentność i rozwiązalność}

\begin{definicja}[Szeregi algebry Liego]
\label{aa:def:lie-series}
Dla algebry Liego $\mathfrak g$ określamy dolny
szereg centralny przez $\gamma_1(\mathfrak g)=\mathfrak g$,
$\gamma_{k+1}(\mathfrak g)=[\mathfrak g,\gamma_k(\mathfrak g)]$
dla $k\ge1$, a szereg pochodny przez
$\mathfrak g^{(0)}=\mathfrak g$,
$\mathfrak g^{(k+1)}=[\mathfrak g^{(k)},\mathfrak g^{(k)}]$
dla $k\ge0$.
Algebra jest nilpotentna, jeśli pewne $\gamma_k=0$,
i rozwiązalna, jeśli pewne $\mathfrak g^{(k)}=0$.
Niezerowa algebra nilpotentna ma klasę $c$, jeśli
$\gamma_c\ne0$ i $\gamma_{c+1}=0$; algebra zerowa
ma klasę zero.
\end{definicja}

\begin{lemat}
\label{aa:lem:lie-series}
Oba szeregi z Definicji~\ref{aa:def:lie-series}
są malejącymi ciągami ideałów. Dla $i,j\ge1$ zachodzi
$[\gamma_i,\gamma_j]\subset\gamma_{i+j}$.
W szczególności $\mathfrak g^{(k)}\subset\gamma_{2^k}$,
więc każda algebra nilpotentna jest rozwiązalna.
\end{lemat}
\begin{proof}
Jeśli $U,W$ są ideałami, to dla $x\in\mathfrak g$,
$u\in U,w\in W$ mamy
$[x,[u,w]]=[[x,u],w]+[u,[x,w]]\in[U,W]$.
Z liniowości wynika, że $[U,W]$ jest ideałem,
zawartym w $U\cap W$. Indukcja daje pierwszą tezę.
Dla ustalonego $i$ dowodzimy drugiej tezy jednocześnie
dla wszystkich $j$ przez indukcję względem $i$.
Dla $i=1$ jest to definicja. Jeśli teza zachodzi
dla $i$ i wszystkich $j$, to z Jacobiego
\[
 [[\mathfrak g,\gamma_i],\gamma_j]
 \subset[\mathfrak g,[\gamma_i,\gamma_j]]
       +[\gamma_i,[\mathfrak g,\gamma_j]]
 \subset\gamma_{i+j+1}.
\]
Zatem zachodzi dla $i+1$. Dla szeregu pochodnego
indukcja od $\mathfrak g^{(0)}=\gamma_1$ daje
$\mathfrak g^{(k+1)}\subset
[\gamma_{2^k},\gamma_{2^k}]\subset\gamma_{2^{k+1}}$.
Jeśli dolny szereg centralny osiąga zero, to przez
malejący charakter ciągu odpowiedni wyraz szeregu
pochodnego jest również zerowy.
\end{proof}

\begin{stwierdzenie}[Podalgebry i obrazy]
\label{aa:prop:lie-series-images}
Podalgebra algebry nilpotentnej (odpowiednio rozwiązalnej)
jest nilpotentna (odpowiednio rozwiązalna).
Jeśli $f:\mathfrak g\to\mathfrak h$ jest surjektywnym
homomorfizmem, to $f(\gamma_k(\mathfrak g))=
\gamma_k(\mathfrak h)$ oraz $f(\mathfrak g^{(k)})=
\mathfrak h^{(k)}$. W szczególności ilorazy przez
ideały dziedziczą obie własności.
\end{stwierdzenie}
\begin{proof}
Dla podalgebry $\mathfrak b$ indukcja daje
$\gamma_k(\mathfrak b)\subset\gamma_k(\mathfrak g)$
oraz $\mathfrak b^{(k)}\subset\mathfrak g^{(k)}$:
nawias zachowuje inkluzje w obu argumentach.
Dla homomorfizmu i podprzestrzeni $U,W$ zachodzi
$f([U,W])=[f(U),f(W)]$, gdyż liniowość $f$
przeprowadza rozpiętość na rozpiętość obrazów,
a nawias na nawias. Stosując to indukcyjnie
i korzystając z surjektywności w pierwszym wyrazie,
otrzymujemy obie równości.
\end{proof}

\begin{przyklad}[Macierze trójkątne]
\label{aa:prz:triangular-lie}
Niech $\mathfrak b_n$ będzie algebrą macierzy
górnotrójkątnych nad $K$, a $\mathfrak n_n$
jej podprzestrzenią macierzy ściśle górnotrójkątnych.
Obie są zamknięte na iloczyny i komutatory.
Dla $r\ge1$ niech $F_r$ będzie rozpiętością
$E_{ij}$ z $j-i\ge r$; dla $r\ge n$ mamy $F_r=0$.
Wzór na iloczyny jednostek macierzowych daje
$F_rF_s\subset F_{r+s}$ (niezerowy iloczyn wymaga
zgodności środkowych indeksów, a różnice indeksów
wówczas się dodają). Zatem $[F_r,F_s]\subset F_{r+s}$.
Mamy $\gamma_1(\mathfrak n_n)=F_1$. Dla $r\ge1$
każda jednostka $E_{ij}$ z $j-i\ge r+1$ spełnia
$E_{ij}=[E_{i,i+1},E_{i+1,j}]$, gdzie drugi
czynnik należy do $F_r$. To daje
$[F_1,F_r]=F_{r+1}$ i indukcyjnie
$\gamma_r(\mathfrak n_n)=F_r$. Stąd dla $n\ge2$
klasa nilpotentności wynosi $n-1$, a dla $n=1$
algebra jest zerowa.

Komutator dwóch macierzy górnotrójkątnych ma zerową
przekątną, bo na $i$-tym wpisie przekątnej iloczynu
pozostaje tylko iloczyn dwóch $i$-tych wpisów.
Ponadto dla $i<j$ mamy $[E_{ii},E_{ij}]=E_{ij}$,
więc $[\mathfrak b_n,\mathfrak b_n]=\mathfrak n_n$.
Z lematu~\ref{aa:lem:lie-series} algebra
$\mathfrak b_n$ jest rozwiązalna.
Dla $n\ge2$ nie jest nilpotentna: skoro
$E_{12}=[E_{11},E_{12}]$, indukcja daje
$E_{12}\in\gamma_k(\mathfrak b_n)$ dla każdego $k$.
\end{przyklad}

\begin{przyklad}[Algebra Heisenberga i algebra $\mathfrak{sl}_2$]
\label{aa:prz:heisenberg-sl2}
Nad ciałem charakterystyki zero algebra
$\mathfrak n_3$ z bazą $x=E_{12},y=E_{23},z=E_{13}$
spełnia $[x,y]=z$, $[x,z]=[y,z]=0$.
Ma centrum $Kz$: dla $u=ax+by+cz$ warunki
$[u,x]=-bz=0$, $[u,y]=az=0$ dają $a=b=0$.
Jej dolny szereg centralny to
$\mathfrak n_3\supset Kz\supset0$.
Dla porównania $\mathfrak{sl}_2$ ma bazę
$h=E_{11}-E_{22}$, $e=E_{12}$, $f=E_{21}$,
z nawiasami $[h,e]=2e$, $[h,f]=-2f$, $[e,f]=h$.
Każdy element tej bazy należy do
$[\mathfrak{sl}_2,\mathfrak{sl}_2]$, więc szereg
pochodny jest stały i algebra nie jest rozwiązalna.
W szczególności nie jest nilpotentna.
\end{przyklad}

\begin{center}
\begin{tikzpicture}[node distance=6mm, every node/.style={font=\small}]
 \node (a) at (0,0) {$\mathfrak n_3$};
 \node (b) at (0,-.9) {$Kz$};
 \node (c) at (0,-1.8) {$0$};
 \draw[-{Stealth}] (a)--node[right] {$[\mathfrak n_3,\cdot]$} (b);
 \draw[-{Stealth}] (b)-- (c);
 \node (d) at (5,0) {$\mathfrak b_2$};
 \node (e) at (5,-.9) {$KE_{12}$};
 \node (f) at (5,-1.8) {$0$};
 \draw[-{Stealth}] (d)--node[right] {$[\mathfrak b_2,\mathfrak b_2]$} (e);
 \draw[-{Stealth}] (e)-- (f);
 \draw[-{Stealth}] (e.west) to[out=170,in=190,looseness=5]
 node[left] {$[\mathfrak b_2,\cdot]$} (e.west);
\end{tikzpicture}

\small Dolny szereg centralny algebry Heisenberga (po lewej).
Szereg pochodny $\mathfrak b_2$ osiąga zero,
ale jej dolny szereg centralny zatrzymuje się na $KE_{12}$.
\end{center}

\section{Twierdzenie Engela}

\begin{lemat}[Komutator z operatorem nilpotentnym]
\label{aa:lem:engel-ad}
Niech $V$ będzie przestrzenią wektorową nad $K$,
a $X\in\operatorname{End}_K(V)$ niech spełnia $X^m=0$
dla pewnego $m\ge1$. Wówczas operator
$\operatorname{ad}_X$ na $\mathfrak{gl}(V)$ jest
nilpotentny, a dokładniej $(\operatorname{ad}_X)^{2m-1}=0$.
\end{lemat}
\begin{proof}
Operatory $L_X(Y)=XY$ i $R_X(Y)=YX$ komutują
z powodu łączności składania. Mamy
$\operatorname{ad}_X=L_X-R_X$.
Ze wzoru dwumianowego dla komutujących operatorów
\[
 (\operatorname{ad}_X)^k(Y)
 =\sum_{j=0}^k(-1)^j\binom{k}{j}X^{k-j}YX^j.
\]
Jeśli $k=2m-1$, to dla każdego $j$ co najmniej
jedna z liczb $j,k-j$ jest nie mniejsza niż $m$.
Każdy składnik znika.
\end{proof}

\begin{twierdzenie}[Twierdzenie Engela dla operatorów]
\label{aa:tw:engel-linear}
Niech $V\ne0$ będzie skończenie wymiarową przestrzenią
wektorową nad dowolnym ciałem $K$, a $\mathfrak l$
podalgebrą $\mathfrak{gl}(V)$. Załóżmy, że każdy
$X\in\mathfrak l$ jest operatorem nilpotentnym na $V$.
Wtedy istnieje $v\in V\setminus\{0\}$ takie, że
$Xv=0$ dla wszystkich $X\in\mathfrak l$.
\end{twierdzenie}
\begin{proof}
Indukcja przebiega względem wymiaru $\mathfrak l$,
jednocześnie dla wszystkich niezerowych skończenie
wymiarowych przestrzeni reprezentacji.
Dla $\mathfrak l=0$ wybieramy dowolny $v\ne0$.
Niech $\dim\mathfrak l>0$ i załóżmy tezę dla
wszystkich mniejszych wymiarów algebr operatorów.

Wybierzmy właściwą podalgebrę $\mathfrak h$ o największym
możliwym wymiarze; istnieje, bo podalgebra zerowa
jest właściwa, a wymiary są liczbami całkowitymi
ograniczonymi przez $\dim\mathfrak l-1$.
Na niezerowej przestrzeni $\mathfrak l/\mathfrak h$
(jedynie ilorazie przestrzeni wektorowych na tym etapie)
określamy działanie
\[
 Y\cdot(X+\mathfrak h)=[Y,X]+\mathfrak h,
 \qquad Y\in\mathfrak h,\ X\in\mathfrak l.
\]
Jest dobrze określone, ponieważ $\mathfrak h$
jest podalgebrą. Tożsamość Jacobiego daje
warunek reprezentacji. Dla $Y\in\mathfrak h$
operator $\operatorname{ad}_Y$ jest nilpotentny
na $\mathfrak{gl}(V)$ na mocy
Lematu~\ref{aa:lem:engel-ad}; jego ograniczenie
do $\mathfrak l$ zachowuje $\mathfrak h$,
więc operator indukowany na ilorazie jest nilpotentny.
Obraz reprezentacji ma wymiar co najwyżej
$\dim\mathfrak h<\dim\mathfrak l$.
Można więc zastosować założenie indukcyjne
do tego obrazu na $\mathfrak l/\mathfrak h$.
Istnieje $X\in\mathfrak l\setminus\mathfrak h$
takie, że $[Y,X]\in\mathfrak h$ dla każdego
$Y\in\mathfrak h$.

Normalizator $N=\{Z\in\mathfrak l:[Z,\mathfrak h]
\subset\mathfrak h\}$ jest podalgebrą:
dla $Z_1,Z_2\in N,Y\in\mathfrak h$ mamy
$[[Z_1,Z_2],Y]=[Z_1,[Z_2,Y]]-[Z_2,[Z_1,Y]]
\in\mathfrak h$, a warunek jest liniowy.
Zawiera $\mathfrak h$ i $X\notin\mathfrak h$,
więc maksymalność wymiaru $\mathfrak h$ wymusza
$N=\mathfrak l$. Oznacza to, że $\mathfrak h$
jest ideałem. Iloraz $\mathfrak l/\mathfrak h$
ma wymiar jeden: gdyby miał wymiar większy,
dowolna prosta w nim byłaby właściwą niezerową
podalgebrą (nawias na prostej jest zerowy),
a jej przeciwobraz byłby właściwą podalgebrą
ściśle zawierającą $\mathfrak h$.

Z założenia indukcyjnego dla $\mathfrak h$ na $V$
przestrzeń $U=\{v\in V:Yv=0\ \forall Y\in\mathfrak h\}$
jest niezerowa. Wybieramy $X\notin\mathfrak h$;
mamy $\mathfrak l=\mathfrak h+KX$.
Dla $u\in U,Y\in\mathfrak h$ zachodzi
$Y(Xu)=X(Yu)+[Y,X]u=0$,
gdyż $Yu=0$ i $[Y,X]\in\mathfrak h$.
Zatem $X$ zachowuje $U$, a jego ograniczenie
do $U$ jest nilpotentne. Ma niezerowe jądro:
gdyby było iniektywne, każda jego potęga byłaby
iniektywna, co przeczy zerowej potędze na $U\ne0$.
Wybieramy $v\ne0$ w tym jądrze. Wówczas
$\mathfrak h v=0$ i $Xv=0$, więc $\mathfrak l v=0$.
\end{proof}

\begin{wniosek}[Wspólna postać ściśle trójkątna]
\label{aa:cor:engel-flag}
Przy założeniach Twierdzenia~\ref{aa:tw:engel-linear}
istnieje pełna flaga podprzestrzeni
$0=V_0\subset V_1\subset\cdots\subset V_d=V$,
$\dim V_i=i$, gdzie $d=\dim V$, taka, że
$X V_i\subset V_{i-1}$ dla każdego $X\in\mathfrak l$
i $1\le i\le d$. W bazie zgodnej z flagą wszystkie
macierze $X$ są ściśle górnotrójkątne.
\end{wniosek}
\begin{proof}
Indukcja po $d$; dla $d=0$ flaga jest trywialna.
Dla $d\ge1$ twierdzenie daje prostą $V_1=Kv$
anihilowaną przez $\mathfrak l$.
Na $V/V_1$ każdy indukowany operator jest nilpotentny,
ponieważ zerowa potęga operatora przechodzi na iloraz.
Obraz $\mathfrak l$ jest podalgebrą operatorów,
więc indukcja stosuje się do ilorazu wymiaru $d-1$.
Przeciwobrazy kolejnych podprzestrzeni flagi ilorazowej
dają $V_1,\ldots,V_d$, a inkluzje obrazów implikują
$XV_i\subset V_{i-1}$ dla $i\ge2$.
Dla $i=1$ wynika to z wyboru $v$.
Wektor bazowy $v_i\in V_i\setminus V_{i-1}$
ma obraz w rozpiętości wcześniejszych wektorów,
co daje żądaną postać macierzy.
\end{proof}

\begin{twierdzenie}[Twierdzenie Engela]
\label{aa:tw:engel}
Niech $\mathfrak g$ będzie skończenie wymiarową
algebrą Liego nad dowolnym ciałem $K$.
Wówczas $\mathfrak g$ jest nilpotentna wtedy
i tylko wtedy, gdy $\operatorname{ad}_x$ jest
operatorem nilpotentnym dla każdego $x\in\mathfrak g$.
\end{twierdzenie}
\begin{proof}
Jeśli $\gamma_{c+1}=0$, to
$\operatorname{ad}_x(\gamma_i)\subset\gamma_{i+1}$,
więc $(\operatorname{ad}_x)^c(\mathfrak g)=0$.
Przypadek $\mathfrak g=0$ jest trywialny.
Odwrotnie, załóżmy nilpotentność wszystkich
$\operatorname{ad}_x$. Na mocy
Stwierdzenia~\ref{aa:prop:adjoint-lie} ich zbiór jest
podalgebrą $\mathfrak{gl}(\mathfrak g)$.
Z Wniosku~\ref{aa:cor:engel-flag} dla $\mathfrak g$
wymiaru $d>0$ otrzymujemy flagę $V_i$ z warunkiem
$[\mathfrak g,V_i]\subset V_{i-1}$.
Ponieważ $\gamma_1=\mathfrak g=V_d$, indukcja daje
$\gamma_j\subset V_{d+1-j}$ dla $1\le j\le d+1$.
W szczególności $\gamma_{d+1}=0$.
\end{proof}

\section{Twierdzenie Liego}

\begin{twierdzenie}[Twierdzenie Liego o wspólnym wektorze własnym]
\label{aa:tw:lie-eigenvector}
Niech $K$ będzie algebraicznie domkniętym ciałem
charakterystyki zero, $\mathfrak g$ skończenie
wymiarową rozwiązalną algebrą Liego nad $K$,
a $\rho:\mathfrak g\to\mathfrak{gl}(V)$ reprezentacją
w niezerowej skończenie wymiarowej przestrzeni $V$.
Wtedy istnieją $v\ne0$ i funkcjonał liniowy
$\lambda:\mathfrak g\to K$ takie, że
$\rho(x)v=\lambda(x)v$ dla każdego $x\in\mathfrak g$.
\end{twierdzenie}
\begin{proof}
Indukcja względem $\dim\mathfrak g$ jednocześnie
dla wszystkich takich reprezentacji. Dla algebry
zerowej bierzemy dowolny wektor niezerowy i funkcjonał
zerowy. Załóżmy $\dim\mathfrak g>0$.
Rozwiązalność wymusza $[\mathfrak g,\mathfrak g]
\ne\mathfrak g$, gdyż inaczej szereg pochodny
byłby stały. Wybieramy hiperpłaszczyznę $\mathfrak h$
zawierającą $[\mathfrak g,\mathfrak g]$ (uzupełniamy
bazę tej właściwej podprzestrzeni do bazy $\mathfrak g$
i pomijamy jeden z nowych wektorów).
Jest ideałem, bo $[\mathfrak g,\mathfrak h]$
zawiera się w $[\mathfrak g,\mathfrak g]$.
Na mocy Stwierdzenia~\ref{aa:prop:lie-series-images}
jest rozwiązalna, więc do ograniczenia reprezentacji
stosujemy założenie indukcyjne.
Istnieją $v\ne0$ i $\mu:\mathfrak h\to K$ liniowe
takie, że $\rho(y)v=\mu(y)v$ dla $y\in\mathfrak h$.

Wybieramy $x\in\mathfrak g\setminus\mathfrak h$
i oznaczamy $T=\rho(x)$. Pokażemy najpierw
$\mu([x,y])=0$ dla wszystkich $y\in\mathfrak h$.
Niech $W_j=\operatorname{span}(v,Tv,\ldots,T^jv)$
dla $j\ge0$, $W_{-1}=0$.
Dla wszystkich $y\in\mathfrak h$ wykażemy indukcyjnie
\[
 \rho(y)T^jv-\mu(y)T^jv\in W_{j-1}.
\]
Dla $j=0$ wynika to z wyboru $v$. Jeśli teza zachodzi
dla wykładnika $j-1$ i wszystkich $y\in\mathfrak h$,
to
\[
 \rho(y)T^jv=T\rho(y)T^{j-1}v+\rho([y,x])T^{j-1}v.
\]
Pierwszy składnik różni się od $\mu(y)T^jv$
o element $TW_{j-2}\subset W_{j-1}$.
Drugi należy do $W_{j-1}$, bo $[y,x]\in\mathfrak h$
i stosujemy to samo założenie indukcyjne dla $[y,x]$.
Daje to żądaną inkluzję.

Niech $r\ge1$ będzie najmniejszą liczbą, dla której
$T^rv$ należy do $\operatorname{span}(v,\ldots,T^{r-1}v)$;
istnieje dzięki skończonemu wymiarowi $V$.
Przestrzeń $W=W_{r-1}$ ma bazę $v,\ldots,T^{r-1}v$
i jest zachowana przez $T$ oraz wszystkie $\rho(y)$
dla $y\in\mathfrak h$ na mocy wykazanej inkluzji.
W tej bazie $\rho(y)|_W$ ma postać górnotrójkątną
z przekątną stale równą $\mu(y)$.
W szczególności dla $[x,y]\in\mathfrak h$
\[
 r\mu([x,y])=\operatorname{tr}(\rho([x,y])|_W)
 =\operatorname{tr}([T|_W,\rho(y)|_W])=0.
\]
Ostatnia równość wynika z cykliczności śladu,
wykazanej w Przykładzie~\ref{aa:prz:classical-lie}.
Charakterystyka zero gwarantuje $r\ne0$ w $K$,
więc $\mu([x,y])=0$.

Rozważmy teraz całą niezerową przestrzeń
$U=\{u\in V:\rho(y)u=\mu(y)u\ \forall y\in\mathfrak h\}$,
zawierającą $v$. Dla $u\in U,y\in\mathfrak h$
mamy
\[
 \rho(y)Tu=T\rho(y)u+\rho([y,x])u
 =\mu(y)Tu+\mu([y,x])u=\mu(y)Tu.
\]
Zatem $T$ zachowuje $U$. Ponieważ $K$ jest
algebraicznie domknięte, wielomian charakterystyczny
$T|_U$ dodatniego stopnia ma pierwiastek $\nu\in K$;
operator $T|_U-\nu I_U$ ma zerowy wyznacznik,
więc istnieje $w\in U\setminus\{0\}$ z $Tw=\nu w$.
Dla $y+ax\in\mathfrak h+Kx=\mathfrak g$ mamy
$\rho(y+ax)w=(\mu(y)+a\nu)w$.
Jednoznaczność zapisu $y+ax$ daje dobrze określony
liniowy funkcjonał $\lambda(y+ax)=\mu(y)+a\nu$.
\end{proof}

\begin{wniosek}[Wspólna postać trójkątna]
\label{aa:cor:lie-flag}
Przy założeniach Twierdzenia~\ref{aa:tw:lie-eigenvector}
istnieje baza $V$, w której wszystkie operatory
$\rho(x)$ są górnotrójkątne. Ponadto
$\rho([\mathfrak g,\mathfrak g])$ składa się
z macierzy ściśle górnotrójkątnych.
\end{wniosek}
\begin{proof}
Indukcja po wymiarze $V$, przy wymiarze zero
nie ma nic do wykazania. Twierdzenie Liego daje
niezmienniczą prostą $Kv$. Na ilorazie $V/Kv$
mamy reprezentację tej samej rozwiązalnej algebry.
Stosując indukcję, otrzymujemy pełną flagę
niezmienniczą na ilorazie. Jej przeciwobrazy
wraz z $0$ tworzą pełną flagę niezmienniczą na $V$.
Baza zgodna z flagą daje postać górnotrójkątną.
Komutatory macierzy górnotrójkątnych mają zerową
przekątną, jak wykazano w
Przykładzie~\ref{aa:prz:triangular-lie}.
Ponieważ $\rho$ zachowuje nawias i sumy liniowe,
teza o całym obrazie algebry pochodnej wynika
z zamkniętości przestrzeni macierzy ściśle
górnotrójkątnych na kombinacje liniowe.
\end{proof}

\begin{wniosek}[Nilpotentność algebry pochodnej]
\label{aa:cor:lie-derived-nilpotent}
Niech $\mathfrak g$ będzie skończenie wymiarową
rozwiązalną algebrą Liego nad algebraicznie domkniętym
ciałem charakterystyki zero. Wówczas ideał
$\mathfrak a=[\mathfrak g,\mathfrak g]$ jest nilpotentny.
\end{wniosek}
\begin{proof}
Stosujemy Wniosek~\ref{aa:cor:lie-flag} do reprezentacji
dołączonej $\operatorname{ad}$ na $\mathfrak g$.
Jeśli $\mathfrak g=0$, teza jest natychmiastowa.
Dla każdego $x\in\mathfrak a$ operator $\operatorname{ad}_x$
jest ściśle górnotrójkątny w pewnej wspólnej bazie,
a zatem nilpotentny: jego $d$-ta potęga jest zerowa,
gdy $d=\dim\mathfrak g$, na mocy własności filtracji
$F_r$ z Przykładu~\ref{aa:prz:triangular-lie}.
Ograniczenie do $\mathfrak a$ jest dobrze określone
(ideał) i nilpotentne. Jest to właśnie operator
dołączony elementu $x$ w algebrze $\mathfrak a$.
Twierdzenie~\ref{aa:tw:engel} daje nilpotentność
$\mathfrak a$.
\end{proof}

\begin{przyklad}[Znaczenie założeń twierdzenia Liego]
\label{aa:prz:lie-hypotheses}
Nad $\mathbb R$ jednowymiarowa abelowa algebra
operatorów generowana przez
$\begin{pmatrix}0&-1\\1&0\end{pmatrix}$
nie ma rzeczywistego niezerowego wektora własnego:
wielomian charakterystyczny $t^2+1$ nie ma pierwiastka
w $\mathbb R$.

W dodatniej charakterystyce $p$ weźmy algebraicznie
domknięte ciało $K$ i $V=K[t]/(t^p)$.
Operator $X$ mnożenia przez klasę $t$ i operator
$Y$ formalnego różniczkowania są dobrze określone:
pochodna $t^pq$ to $t^pq'$ w charakterystyce $p$.
Na klasie dowolnego wielomianu mamy
$[Y,X]v=Y(tv)-tY(v)=v$, więc $[Y,X]=I$.
Przestrzeń $\mathfrak l=\operatorname{span}(X,Y,I)$
jest algebrą Liego o algebrze pochodnej $KI$
i kolejnym wyrazie zerowym, a więc jest rozwiązalna.
Gdyby $v\ne0$ był wspólnym wektorem własnym $X,Y$,
to $[Y,X]v=0$ (iloczyny skalarów komutują),
co przeczyłoby $[Y,X]v=v$.
Brak zatem wspólnego wektora własnego, mimo
algebraicznej domkniętości ciała.
\end{przyklad}

Twierdzenie Engela dotyczy nilpotentności operatorów,
twierdzenie Liego dotyczy rozwiązalności algebry.
Nie należy utożsamiać tych założeń: algebra abelowa
generowana przez $I$ jest nilpotentna jako algebra
Liego, ale $I$ na niezerowej przestrzeni nie jest
operatorem nilpotentnym. Wniosek o ściśle trójkątnej
postaci z twierdzenia Engela wymaga nilpotentności
każdego operatora w rozważanej reprezentacji.
Kryteria półprostoty i związek z układami korzeni
będą przedmiotem dalszego rozdziału.
